\documentclass[14pt,a4paper]{article}
\usepackage[utf8]{inputenc}
\usepackage[paper=a4paper, margin=1.5cm]{geometry}

% FA icons
\usepackage{hyperref}
\usepackage{url}
\usepackage{color}
\usepackage{graphicx}
\hypersetup{
  colorlinks=false,% hyperlinks will be black
  pdfborderstyle={/S/U/W .5}% border style will be underline of width 1pt
}

% citace a jazyk
\usepackage[main=czech, english]{babel}
\usepackage{csquotes}
\usepackage[style=iso-authoryear]{biblatex}
\addbibresource{sources.bib}

\begin{document}
	\begin{titlepage}
		\begin{center}
            {
            \centering
            \includegraphics[]{./PDF_tisk/UP_logo_PdF-UP_horizont_cz.pdf}
            }
			
			\vspace{3cm}

            {
                \LARGE
                \textbf{Šimon Janča}\\
                3.\,ročník\\[8mm]
                Obor: Informační technologie a Matematika pro vzdělávání
            }

            \vspace{4cm}
			
			{
			    \textbf{\Huge Využití IT ke generování zadání písemných testů}\\[4mm]
			    \Large
			    Bakalářská práce
			}

            \vfill
            
            {
                Vedoucí práce:
                doc. RNDr. Petr Šaloun, Ph.D.
                \hfill
    			Olomouc \the\year{}
            }
			
		\end{center}
	\end{titlepage}
	% generate table of contents
	\tableofcontents
	
	\newpage
	
	\section{Úvod}
    V~této práci si~kladu za~cíl vytvořit webovou aplikaci, kterou by~učitelé mohli používat k~tvorbě různých variant zadání písemných prací.
    Zároveň bych rád popsal základní prvky vývoje moderní \emph{MVC} aplikace a~popsal důvody k~jednotlivým rozhodnutím od~návrhu
    a~výběru technologií až k~její implementaci za~použití vybraných technologií.
    
    Pochopitelně nejde o~jediný způsob, jak aplikace mohou vzniknout, nicméně by~práce měla být čitelným a~reprodukovatelným postupem vývoje
    a~údržby podobné aplikace.
        
        \subsection{V současnosti používané nástroje na generování prací}
        % TODO: 
        \subsection{Existující nástroje}

        \subsection{Účel a~popis fungování aplikace}
        Aplikace \emph{Písemkomat}, poskytne registrovaným uživatelům možnost vytvořit vytvořit si příklady, které může umístit do~volitelně
        strukturovaných kategorií - podle předmětu, ročníku, tématu... Tyto příklady může následně využít při tvorbě písemných prací,
        kde si bude moci zvolit, kategorie ze kterých se příklady vyberou a~počet příkladů, které se~mají do~písemky vygenerovat. 
        \subsection{Příklad použití}
        Mějme prvního uživatele, říkejme mu~třeba Adam. Adam je~učitel a~chce vytvořit 2~verze písemky s~rozdílnými příklady pro~své žáky.
        Adam se~zaregistruje do~aplikace a~přihlásí se~do~svého účtu. Adam učí Matematiku, takže si~vytvoří 5 příkladů z~tohoto předmětu
        a~vloží je do kategorie Matematika.

        Následně si vytvoří písemnou práci a díky tomu, že si~Adam vytvořil 5~příkladů, může si~vybrat, zda chce,
        aby se~v~jeho písemce objevily všechny příklady, nebo jen některé. Bude tedy požadovat 3 příklady s~náhodně dosazenými hodnotami.
        Některé hodnoty nemohou být zcela náhodné, protože by~to mohlo vést k~neřešitelným příkladům, nebo příkladům s~nulovým výsledkem.
        Proto aplikace umožní vytvořit si~proměnné, které budou mít přiřazenou hodnotu z~určitého intervalu nebo vypočtené z~jiných proměnných.

        Adam si~např.~vytvoří proměnné $A, B, C$ a $D$, kde $A, B$ budou náhodná čísla od 1 do 10, $C = 4/3*A$ a $D = C + B$.

        Řešený příklad se~vygeneruje do vzorce, který posléze zapíše do~\emph{Písemkomatu}. Do proměnných $A, B$ se uloží náhodné číslo
        z~intervalu, který~si~Adam nastavil a~pomocí zápisu proměnné - \texttt{\$A} nebo formátu umožňující
        delší názvy\,--\,\texttt{\$\{nazevpromenne\}}\,--\,se~vloží do~příkladu na~zadané místo, resp.~do zápisu výpočtu dalších proměnných,
        kde budou použity ve výpočtu. Tedy v~tomto případě předpokládáme, že~budou vygenerována náhodná čísla $A = 3$ a~$B = 5$
        a~hodnoty proměnných $C, D$ se~dopočítají takto:
        \begin{align*}
            A &= 3 \\
            B &= 5 \\
            C &= 4/3*A = 4/3*3 = 4 \\
            D &= C + B = 4 + 5 = 9
        \end{align*}

        Proměnné $A, B$ se při generování výsledné podoby příkladu nahradí náhodnými čísly z~určeného intervalu, proměnné $C, D$ se vyčíslí.
        Výsledná podoba příkladu může vypadat např.~takto:
        \begin{align*}
            \$Ax + \$B &= \$C + \$D \\
            3x + 5 &= 4 + 9
        \end{align*}

        Adam bude mít taky možnost zvolit~si, jestli chce otevřenou nebo uzavřenou formu příkladu s možnostmi výběru.
        V~tomto případě může taky využít proměnné, které si vytvořil dříve. V obou případech může taky zapsat vzorec,
        ze~kterého se~vypočtou možnosti a označit správnou možnost.
        Ve výsledném příkladu se~v~možnostech výběru proměnné nahradí a~výraz se~vyčíslí.
        Výsledná podoba zadání a~možností příkladu může vypadat např.~takto:
        
        $$ 3x + 5 &= 4 + 9 $$
        \begin{itemize} %alfa
            \item $x = 2$
            \item $x = 3$
            \item $x = 4$
            \item $x = 5$
        \end{itemize}
        
    \section{Vývoj moderních aplikací}
    Pokud se~chceme pustit do~vývoje webové aplikace, která bude využívat výhod moderních postupů a~technologií, je~potřeba
    si~ujasnit pár\,--\,i~historických pojmů a~pochopit, jak~spolu souvisí.
        \subsection{Základní statická HTML stránka}
        \subsection{CSS Styly}
        \subsection{Javascript}
        \subsection{HTML s dynamickým obsahem}
        \subsection{Client-Server}
        \subsection{Kompilátor a Interpret}
        \subsection{MVC - Model View Controller}
        \subsection{Verzovací systém GIT}
        \subsection{Linux a příkazová řádka, WSL2}
        \subsection{JSON - Javascript Object Notation}
        \subsection{API - Application Interface}
        \subsection{AJAX - Asynchronous Javascript Request}
        \subsection{REST - Representational state transfer}
        \subsection{Frontend, Backend}
        \subsection{Webová služba, webový server}
        \subsection{Port}
        \subsection{Zabezpečení}
            \subsubsection{Autentikace}
            \subsubsection{Autorizace}
            \subsubsection{JWT - JSON Web Token}
        \subsection{Údržba aplikace}
            \subsubsection{Sledování bezpečnostních trendů}
            \subsubsection{Sledování chyb nahlášených uživateli}
            \subsubsection{Sentry}
	
	\section{Návrh aplikace}
        \subsection{UCD - Use Case Diagram}
        \subsection{ERD - Entity Relationship Diagram}
        \subsection{Návrh struktury databáze}
        \subsection{Výběr technologií}
            \subsubsection{React, Vite, Typescript}
            \subsubsection{Go}
            \subsubsection{Gin}
            \subsubsection{GORM}
            \subsubsection{MySQL}
            \subsubsection{Docker, docker-compose}
            \subsubsection{Nginx}
	
	\section{Postup implementace}
        \subsection{Backend}
        \subsection{Frontend}
        \subsection{Testování}
        \subsection{Údržba}
        
	\section{Budoucí rozvoj aplikace}
        \subsection{Integrace do jiných prostředí}
        \subsection{Režim cvičení}
        \subsection{Výukové materiály a minikvízy}
        \subsection{Online zkoušení a výuka}
	
	\section{Závěr}
	
	
%\parencite{knuth1990virtual}

\printbibliography
\end{document}
